\PassOptionsToPackage{table,xcdraw}{xcolor}
   \usepackage[
	top = 3 cm, 
	right = 2 cm, 
	left = 3.5 cm, 
	bottom = 3 cm
]{geometry}

\usepackage{iftex}
\ifPDFTeX
    \usepackage[utf8]{inputenc}
    \IfFileExists{t5enc.def}{
        \usepackage[T5]{fontenc}
        \IfFileExists{urwvn.map}{\pdfmapfile{+urwvn.map}}{}
    }{
        \PackageError{report}
        {Missing Vietnamese TeX support: t5enc.def}
        {Install the Vietnamese/other-language TeX package, e.g. `sudo apt install texlive-lang-other`, then rebuild with latexmk.}
    }
\else
    \usepackage{fontspec}
    \setmainfont{TeX Gyre Termes}
    \setsansfont[Scale=0.92]{TeX Gyre Heros}
    \setmonofont{Latin Modern Mono}
\fi
\usepackage[main=vietnamese,provide=*]{babel}
\usepackage{textcomp} 
\ifPDFTeX
    \usepackage[scaled=0.92]{helvet}
    \renewcommand{\rmdefault}{ptm}
\fi

\usepackage{cite}
\usepackage{bibentry}
\usepackage{url}
\usepackage{hyperref}


\usepackage{upquote}
\usepackage{amsmath}
\usepackage[capitalise,nameinlink,noabbrev]{cleveref}
\crefname{chapter}{Chương}{Chương}
\crefname{section}{Mục}{Mục}
\crefname{subsection}{Mục}{Mục}
\crefname{figure}{Hình}{Hình}
\crefname{table}{Bảng}{Bảng}
\Crefname{chapter}{Chương}{Chương}
\Crefname{section}{Mục}{Mục}
\Crefname{subsection}{Mục}{Mục}
\Crefname{figure}{Hình}{Hình}
\Crefname{table}{Bảng}{Bảng}
\usepackage{tikz}
\usetikzlibrary{decorations.pathreplacing}
\usepackage{algorithm}
\floatname{algorithm}{Thuật toán}
\usepackage{booktabs} 
\usepackage{makecell}

\usepackage{listings}


\renewcommand{\listalgorithmname}{Danh sách thuật toán}
\usepackage{algpseudocode}	
\usepackage{amssymb}
\IfFileExists{fontawesome5.sty}{
    \usepackage{fontawesome5}
}{
    \newcommand{\faFolderOpen}{\textbf{[Folder]}}
}
\usepackage{enumitem}
\usepackage{graphicx}
\ifPDFTeX
    \usepackage{mathptmx}
\else
    \usepackage{unicode-math}
    \setmathfont{TeX Gyre Termes Math}
\fi
\usepackage{xurl}
\usepackage{rotating}
\usepackage{tcolorbox}
\tcbuselibrary{listings,skins,breakable}

\usepackage{listings}
\usepackage{caption}
\captionsetup{width=\textwidth}
\usepackage{longtable}
\usepackage{xcolor}
\usepackage{float}
\usepackage{array}
\newcolumntype{L}[1]{>{\raggedright\arraybackslash}p{#1}}
\setlength{\emergencystretch}{2em}

\setlist[enumerate]{itemsep=0pt,topsep=0pt}
\setlist[itemize]{itemsep=0pt,topsep=0pt}
\usepackage{subcaption}
% Thiết lập khoảng cách giữa môi trường enumerate và các đoạn khác
% \setlength{\topsep}{0pt}
% \setlength{\partopsep}{0pt}

% Thiết lập khoảng cách giữa môi trường itemize và các đoạn khác
% \setlength{\topsep}{0pt plus 1pt minus 1pt}
% \setlength{\partopsep}{0pt plus 1pt minus 1pt}


\usepackage{parskip}
\setlength{\parskip}{5pt} 
\setlength{\parindent}{1.25cm}

\usepackage{listings}
\lstset{
    language=Python,
    basicstyle=\ttfamily,
    keywordstyle=\color{blue},
    stringstyle=\color{orange},
    commentstyle=\color{green!70!black},
    morecomment=[l][\color{magenta}]{\#},
    frame=single,
    breaklines=true,
    numbers=left,
    numberstyle=\tiny\color{gray},
    stepnumber=1,
    tabsize=4,
    showstringspaces=false,
    captionpos=b,
    xleftmargin=\parindent,
}

\makeatletter
\g@addto@macro\normalsize{%
  \setlength\abovedisplayskip     {0pt}
  \setlength\belowdisplayskip     {0pt}
  \setlength\abovedisplayshortskip{0pt}
  \setlength\belowdisplayshortskip{0pt}
}
\makeatother

\setlength{\jot}{0pt}

\numberwithin{equation}{chapter}

\usepackage{xcolor}
\definecolor{maincolor}{RGB}{0, 105, 166}
\definecolor{lightgray}{RGB}{242, 242, 242}
\definecolor{olivegreen}{rgb}{0.33, 0.42, 0.18}


\usepackage{tabu}

\usepackage{framed}
\definecolor{shadecolor}{named}{lightgray}

\newcommand{\maincolorit}[1]{{\color{maincolor}\textit{#1}}}
\newcommand{\maincolorbf}[1]{{\color{maincolor}\textbf{#1}}}
\newcommand{\maincolor}  [1]{{\color{maincolor} #1}}
\newcommand{\red}        [1]{{\color{red} #1}}

\usepackage{tikz}
\usepackage{tikz-3dplot}
\usetikzlibrary{calc,math}
\usepackage{tasks}

\usepackage{thmtools}

\tcbuselibrary{most}
\tcbset{
	customstyle/.style={
  		breakable,
  		enhanced,
  		outer arc = 4pt,
  		arc       = 4pt,
  		%colframe  = maincolor,
  		%colback   = maincolor!12,
		colback   = gray!10!white,
  		attach boxed title to top left,
		drop fuzzy shadow,
  		boxed title style = {
  		  	colback   = maincolor,
  		  	outer arc = 4pt,
  		  	arc       = 0pt,
  		  	top       = 1pt,
  		  	bottom    = 1pt,
  		},
  		fonttitle = \bfseries
  	}
}

% Định nghĩa tcolorbox mà không có tiêu đề mặc định
\newtcolorbox{tomlquotebox}[1][]{
    customstyle,
    colframe=maincolor,
    fonttitle=\bfseries,
    #1
}

% Sử dụng gói lệnh minted để định dạng đoạn code
\usepackage{minted}

% Định nghĩa một môi trường để chứa đoạn văn bản TOML
\newenvironment{tomlcode}
{\minted[
    breaklines=true,
    style=tango,
    linenos=false,
    fontsize=\footnotesize
]{toml}}
{\endminted}

\newtcolorbox[auto counter, number within = chapter]{theorembox}[1][]{
  	customstyle,
  	title = Tập thuộc tính ,
  	overlay unbroken and first = {
      	\path
        let
        \p1 = (title.north east),
        \p2 = (frame.north east)
        in
        node[
			anchor = west, 
			color  = maincolor, 
			text width = \x2 - \x1
		] 
        at (title.east) {\faFolderOpen~\textit{#1}}; %Ký hiệu ở đây
    }
}
\colorlet{tcbcol@back}{tcbcolback}
\colorlet{tcbcol@frame}{tcbcolframe}
\newtcolorbox{note}[1][]{
    enhanced,
	before skip = 2mm,
    after skip  = 3mm,
	boxrule     = 0pt,
    left        = 5mm,
    right       = 2mm,
    top         = 1mm,
    bottom      = 1mm,
	colback     = gray!10!white,
	colframe    = white,
	sharp corners, 
    rounded corners = southeast,
    arc is angular,
    arc         = 3mm,
	underlay = {%
		\path[fill = tcbcol@back!80!white] ([yshift=3mm]interior.south east)--++(-0.4,-0.1)--++(0.1,-0.2);
		\path[draw = tcbcol@frame, shorten <= -0.05mm, shorten >= -0.05mm] ([yshift=3mm]interior.south east)--++(-0.4,-0.1)--++(0.1,-0.2);
		\path[fill = maincolor, draw = none] (interior.south west) rectangle node[white]{\huge\bfseries !}
		([xshift=4mm]interior.north west);
	},
	drop fuzzy shadow, #1}

\newtcolorbox[auto counter, number within = chapter]{examplebox}{
	breakable,
	enhanced,
	colback      = gray!10!white,
	boxrule      = 0pt,
	arc          = 0pt,
	outer arc    = 0pt,
	title        = Ví dụ~\thetcbcounter,
	fonttitle    = \bfseries,
	coltitle     = white,
	colbacktitle = white,
	colframe     = white,
	title style  = maincolor,
	drop fuzzy shadow,
	overlay = {
	  	\draw[line width = 2pt, maincolor] 
			(frame.south west) -- (frame.south east);
	}
}
\usepackage{footnote}
\usepackage{fancyhdr}
\pagestyle{fancy}
\fancyhf{}
\fancyfoot[CO]{
	\noindent 
	\tikz[baseline]{
		\draw[line width = 0.4mm, maincolor] (-10,0) -- (5,0);
		\draw[fill = maincolor, draw = maincolor] 
			(5-0.8,0)--(5.3-0.8,0.3)--(6.5-0.8,0.3)--(6.2-0.8,0)--(6.5-0.8,-0.3)--(5.3-0.8,-0.3)--cycle; 
		\draw[fill = maincolor!50!white, draw = maincolor!50!white] 
			(6.35-0.8,0)--(6.65-0.8,0.3)--(6.85-0.8,0.3)--(6.55-0.8,0)--(6.85-0.8,-0.3)--(6.65-0.8,-0.3)--cycle;
		\draw[fill = maincolor, draw = maincolor]
			(6.7-0.8,0)--(7-0.8,0.3)--(9.5-0.8,0.3)--(9.5-0.8,-0.3)--(7-0.8,-0.3)--cycle; 
		\node at (5.2-0.8,0.02)[anchor = west]{\color{white}\textbf{\thepage}};
		%\node at (-10,-0.5)[anchor = west]{\color{white}Mẫu \LaTeX\, này được thiết kế bởi Lê Nguyên Bách 20195841};
	}
}
\fancyfoot[CE]{
	\noindent 
	\tikz[baseline, overlay, remember picture]{
		%\node at (15,0) {\,};
		\draw[line width = 0.4mm, maincolor] (-6.2,0) -- (8,0);
		\draw[fill = maincolor, draw = maincolor] 
			(-5-1.2,0)--(-5.3-1.2,0.3)--(-6.5-1.2,0.3)--(-6.2-1.2,0)--(-6.5-1.2,-0.3)--(-5.3-1.2,-0.3)--cycle;
		\draw[fill = maincolor!50!white, draw = maincolor!50!white] 
			(-6.35-1.2,0)--(-6.65-1.2,0.3)--(-6.85-1.2,0.3)--(-6.55-1.2,0)--(-6.85-1.2,-0.3)--(-6.65-1.2,-0.3)--cycle;
		\draw[fill = maincolor, draw = maincolor]
			(-6.7-1.2,0)--(-7-1.2,0.3)--(-9.5-1.2,0.3)--(-9.5-1.2,-0.3)--(-7-1.2,-0.3)--cycle; 
		\node at (-6.4,0.02)[anchor = east]{\color{white}\textbf{\thepage}};
		%\node at (8.8,-0.5) [anchor = east]{\color{white}Mẫu \LaTeX\, này được thiết kế bởi Lê Nguyên Bách 20195841};
	}
}

\fancypagestyle{plain}{%
\fancyfoot[CO]{
	\noindent 
	\tikz[baseline]{
		\draw[line width = 0.4mm, maincolor] (-10,0) -- (5,0);
		\draw[fill = maincolor, draw = maincolor] 
			(5-0.8,0)--(5.3-0.8,0.3)--(6.5-0.8,0.3)--(6.2-0.8,0)--(6.5-0.8,-0.3)--(5.3-0.8,-0.3)--cycle; 
		\draw[fill = maincolor!50!white, draw = maincolor!50!white] 
			(6.35-0.8,0)--(6.65-0.8,0.3)--(6.85-0.8,0.3)--(6.55-0.8,0)--(6.85-0.8,-0.3)--(6.65-0.8,-0.3)--cycle;
		\draw[fill = maincolor, draw = maincolor]
			(6.7-0.8,0)--(7-0.8,0.3)--(9.5-0.8,0.3)--(9.5-0.8,-0.3)--(7-0.8,-0.3)--cycle; 
		\node at (5.2-0.8,0.02)[anchor = west]{\color{white}\textbf{\thepage}};
		%\node at (-10,-0.5)[anchor = west]{\color{white}Mẫu \LaTeX\, này được thiết kế bởi Lê Nguyên Bách 20195841};
	}
}
\fancyfoot[CE]{
	\noindent 
	\tikz[baseline, overlay, remember picture]{
		%\node at (15,0) {\,};
		\draw[line width = 0.4mm, maincolor] (-6.2,0) -- (8,0);
		\draw[fill = maincolor, draw = maincolor] 
			(-5-1.2,0)--(-5.3-1.2,0.3)--(-6.5-1.2,0.3)--(-6.2-1.2,0)--(-6.5-1.2,-0.3)--(-5.3-1.2,-0.3)--cycle;
		\draw[fill = maincolor!50!white, draw = maincolor!50!white] 
			(-6.35-1.2,0)--(-6.65-1.2,0.3)--(-6.85-1.2,0.3)--(-6.55-1.2,0)--(-6.85-1.2,-0.3)--(-6.65-1.2,-0.3)--cycle;
		\draw[fill = maincolor, draw = maincolor]
			(-6.7-1.2,0)--(-7-1.2,0.3)--(-9.5-1.2,0.3)--(-9.5-1.2,-0.3)--(-7-1.2,-0.3)--cycle; 
		\node at (-6.4,0.02)[anchor = east]{\color{white}\textbf{\thepage}};
		%\node at (8.8,-0.5) [anchor = east]{\color{white}Mẫu \LaTeX\, này được thiết kế bởi Lê Nguyên Bách 20195841};
	}
}
}

\renewcommand{\headrulewidth}{0pt}

\usepackage{hyperref}
\hypersetup{
    colorlinks  = true,
    citecolor   = maincolor!90,
    linkcolor   = maincolor,
    filecolor   = magenta,      
    urlcolor    = maincolor,
}

\renewcommand{\labelitemi}{$\textcolor{maincolor}{\bullet}$}
\renewcommand{\labelitemii}{$\textcolor{maincolor}{\cdot}$}
\renewcommand{\labelitemiii}{$\textcolor{maincolor}{\diamond}$}
\renewcommand{\labelitemiv}{$\textcolor{maincolor}{\ast}$}


\usepackage[explicit]{titlesec}
\titleformat{\chapter}[display]
    {\normalfont\Large\rmfamily}
    {\flushright\fontsize{60}{0}\textbf{\textcolor{maincolor}{{\huge\chaptername}~\thechapter\vskip0pt\rule{\textwidth}{2pt}}}}{0pt}
    {\flushleft\fontsize{30}{0}{\maincolor{\textbf{#1}}}\vskip 10pt}
\titlespacing*{\chapter}
    {0pt}{-40pt}{0pt}


\titleformat{\section}
{\color{maincolor}\normalfont\Large\bfseries}
{\color{maincolor}\thesection}{5pt}{\color{maincolor}#1}

% Subsection title style 
\titleformat{\subsection}
{\color{maincolor}\normalfont\large\bfseries}
{\color{maincolor}\thesubsection}{5pt}{\color{maincolor}#1}

% Subsubsection title style 
\titleformat{\subsubsection}
{\color{maincolor}\normalfont\normalsize\bfseries}
{\color{maincolor}\thesubsubsection}{5pt}{\color{maincolor}#1}
